\documentclass[../libro.tex]{subfiles}

\begin{document}

\ifSubfilesClassLoaded{\mainmatter\chapter{行列式}\clearpage}{}

\KunAsteriskoEnEnhavtabelo
\section{\texorpdfstring{由 \(m\)~个 \(n\)~元
      \({\leq} 1\)~次方程作成的方程组 (1)}%
  {由 m 个 n 元 ≤1 次方程作成的方程组 (1)}}
\SenAsteriskoEnEnhavtabelo

\maldevigalegajxo

在接下来的若干节里, 我想用行列式讨论%
由 \(m\)~个 \(n\)~元 \({\leq} 1\)~次方程作成的方程组%
的解.
注意, 方程数 \(m\) 不一定等于未知数数 \(n\).

我们有时会用 Cramer 公式辅助讨论.
具体地, 为讨论一般的%
由 \(m\)~个 \(n\)~元 \({\leq} 1\)~次方程作成的方程组%
的解,
我们有时会作辅助的线性方程组,
其的方程数与未知数数相等,
且 Cramer 公式是可用的.

本节, 我们讨论, 线性方程组有解时, 解是否唯一.

我们先看一个简单的事实.
注意, 对任何 \(m \times n\)~阵 \(A\),
\(A 0 = 0\),
其中, 左侧的 \(0\) 是 \(n \times 1\)~零阵,
而右侧的 \(0\) 是 \(m \times 1\)~零阵.
故 \(A X = 0\) 必有零解.

\begin{theorem}
    设 \(A\) 是 \(m \times n\)~阵.
    设 \(B\) 是 \(m \times 1\)~阵.
    设 \(n \times 1\)~阵 \(C\) 适合 \(AC = B\).

    (1)
    若 \(AX = 0\) 只有零解, 则 \(AX = B\) 的解唯一;

    (2)
    若存在非零的 \(n \times 1\)~阵 \(D\)
    使 \(AD = 0\),
    则 \(AX = B\) 的解不唯一.
\end{theorem}

\begin{proof}
    (1)
    设 \(n \times 1\)~阵 \(Y\) 适合 \(AY = B\).
    则
    \begin{align*}
        A(Y - C) = AY - AC = B - B = 0.
    \end{align*}
    故 \(Y - C\) 是 \(AX = 0\) 的一个解.
    因为 \(AX = 0\) 只有零解,
    故 \(Y - C = 0\).
    从而 \(Y = C\).

    % The following proof assumes that the underlying field
    % has infinitely many elements.

    % (2)
    % 要证 \(AX = B\) 有无限多个解,
    % 只要证,
    % 对任何正整数 \(\ell\),
    % 我们能找到 \(AX = B\) 的
    % \(\ell\) 个互不相同的解.

    % 作 \(\ell\)~个 \(n \times 1\)~阵
    % \(C_1\), \(C_2\), \(\dots\), \(C_\ell\),
    % 其中, \(C_j = C + jD\).
    % 首先, 每个 \(C_j\) 都是 \(AX = B\) 的解:
    % \begin{align*}
    %     AC_j = A(C + jD) = AC + A(jD) = B + j(AD) = B + j0 = B.
    % \end{align*}
    % 然后, 我们证,
    % 若 \(p\), \(q\) 是二个不相等的%
    % 且不超过 \(\ell\) 的正整数,
    % 则 \(C_p \neq C_q\).
    % 用反证法.
    % 反设 \(C_p = C_q\),
    % 则
    % \begin{align*}
    %     0 = C_p - C_q = (C + pD) - (C + qD) = (p - q)D.
    % \end{align*}
    % 因为 \(p \neq q\), 故 \(p - q \neq 0\).
    % 从而
    % \begin{align*}
    %     0 = \frac{1}{p-q}\, 0
    %     = \frac{1}{p-q}\, ((p-q)D)
    %     = \Big( \frac{1}{p-q}\, (p-q) \Big) D
    %     = 1D
    %     = D.
    % \end{align*}
    % 这是矛盾.
    % 所以, \(C_p \neq C_q\).

    (2)
    因为 \(D \neq 0\), 故 \(C + D \neq C\).
    因为 \(A(C + D) = AC + AD = B + 0 = B\),
    且 \(C + D \neq C\),
    故 \(AX = B\) 的解不唯一.
\end{proof}

由此可见, 若 \(AX = 0\) 有非零解,
则 \(AX = B\) 有解时,
它的解不唯一.
若 \(AX = 0\) 只有零解,
则 \(AX = B\) 有解时,
它的解唯一.

然后, 我们讨论, \(AX = 0\) 是否有非零解.
不过, 为使讨论简单, 我们要作准备.

\begin{theorem}
    设 \(A\) 是 \(m \times n\)~阵.
    若 \(A\) 没有行列式非零的 \(r+1\)~级子阵,
    则对任何高于 \(r\)~的整数 \(s\),
    \(A\) 没有行列式非零的 \(s\)~级子阵.
\end{theorem}

说阵~\(Z\) 是 \(A\) 的一个 \emph{\(p\)~级子阵},
就是说, \(Z\) 是 \(A\) 的一个子阵
(见本章, 节~\sekcio{5}),
且 \(Z\) 是一个 \(p\)~级阵.

\begin{proof}
    我们用数学归纳法证明此事.
    具体地, 设命题 \(P(s)\) 为:
    \begin{quotation}
        \(A\) 没有行列式非零的 \(s\)~级子阵.
    \end{quotation}
    则我们的目标是,
    对任何高于 \(r\)~的整数 \(s\),
    \(P(s)\) 是对的.

    \(P(r+1)\) 是对的.

    假定 \(P(t)\) 是对的.
    我们由此证, \(P(t+1)\) 也是对的.
    若 \(A\) 没有 \(t+1\)~级子阵,
    此事自然是对的.
    若 \(A\) 有 \(t+1\)~级子阵,
    按列~\(1\) 展开它的行列式,
    可知,
    此子阵的行列式%
    是 \(A\)~的 \(t+1\)~个 \(t\)~级子阵的行列式的倍的和.
    由此可知,
    \(A\) 没有行列式非零的 \(t+1\)~级子阵.

    所以, \(P(t+1)\) 是对的.
    由数学归纳法, 待证命题成立.
\end{proof}

\begin{theorem}
    设 \(A\) 是 \(m \times n\)~阵.
    存在一个唯一的非负整数 \(r\),
    使:

    (1)
    \(A\) 有一个行列式非零的 \(r\)~级子阵;

    (2)
    \(A\) 没有行列式非零的 \(r+1\)~级子阵.
\end{theorem}

\begin{proof}
    唯一性是简单的.
    设还有非负整数 \(s\) 适合条件:

    (1\('\))
    \(A\) 有一个行列式非零的 \(s\)~级子阵;

    (2\('\))
    \(A\) 没有行列式非零的 \(s+1\)~级子阵.

    反设 \(s > r\).
    既然 \(A\) 没有行列式非零的 \(r+1\)~级子阵,
    故 \(A\) 没有行列式非零的 \(s\)~级子阵.
    这跟 (1\('\)) 矛盾.
    所以, \(s \leq r\).
    类似地, 我们可证 \(r \leq s\).
    从而 \(s = r\).

    下面, 我们说明存在性.
    我们设 \(t\) 是 \(m\) 与 \(n\) 中的较小者.
    % 若此 \(r\) 存在, 则它既不超过 \(m\), 也不超过 \(n\).
    若 \(A\) 有行列式非零的 \(t\)~级子阵,
    我们可取 \(r = t\);
    否则, 我们代 \(t\) 以 \(t - 1\).
    若 \(A\) 有行列式非零的 \(t - 1\)~级子阵,
    我们可取 \(r = t - 1\);
    否则, 我们代 \(t - 1\) 以 \(t - 2\).
    ……
    反复地作下去,
    我们能找到此 \(r\).
    (注意, 我们约定 ``\(0\)~级阵'' 的行列式为 \(1\),
    所以这个过程能结束.)
\end{proof}

\begin{theorem}[\(0\) 的作用]
    设 \(A\) 是 \(m \times n\)~阵.
    设 \(B\) 是 \(m \times 1\)~阵.
    设 \(k\) 是一个不超过 \(m\)~的正整数.
    作一个 \((m + 1) \times n\)~阵 \(\underbar{A}\),
    其中,
    \begin{align*}
        [\underbar{A}]_{i,j}
        = \begin{cases}
              [A]_{i,j},   & i \leq k;  \\
              0,           & i = k + 1; \\
              [A]_{i-1,j}, & i > k + 1.
          \end{cases}
    \end{align*}
    (通俗地, \(\underbar{A}\)
    是在 \(A\)~的行~\(k\) 下写一行 \(0\) 作成的阵.)
    再作一个 \((m + 1) \times 1\)~阵 \(\underbar{B}\),
    其中,
    \begin{align*}
        [\underbar{B}]_{i,1}
        = \begin{cases}
              [B]_{i,1},   & i \leq k;  \\
              0,           & i = k + 1; \\
              [B]_{i-1,1}, & i > k + 1.
          \end{cases}
    \end{align*}
    (通俗地, \(\underbar{B}\)
    是在 \(B\)~的行~\(k\) 下写 \(0\) 作成的阵.)

    (1)
    若 \(n \times 1\)~阵 \(C\) 适合
    \(AC = B\),
    则 \(\underbar{A}C = \underbar{B}\).

    (2)
    若 \(n \times 1\)~阵 \(C\) 适合
    \(\underbar{A}C = \underbar{B}\),
    则 \(AC = B\).

    (3)
    若 \(A\) 有一个行列式非零的 \(r\)~级子阵,
    但没有行列式非零的 \(r+1\)~级子阵,
    则 \(\underbar{A}\) 有一个行列式非零的 \(r\)~级子阵,
    但没有行列式非零的 \(r+1\)~级子阵.

    (4)
    若 \(\underbar{A}\) 有一个行列式非零的 \(r\)~级子阵,
    但没有行列式非零的 \(r+1\)~级子阵,
    则 \(A\) 有一个行列式非零的 \(r\)~级子阵,
    但没有行列式非零的 \(r+1\)~级子阵.
\end{theorem}

\begin{proof}
    (1)
    设 \(AC = B\).
    我们要证 \(\underbar{A}C = \underbar{B}\).

    首先, \(\underbar{A}C\) 与 \(\underbar{B}\)
    的尺寸都是 \((m+1) \times 1\).

    若 \(i \leq k\), 则
    \begin{align*}
        [\underbar{A}C]_{i,1}
        = {} &
        \sum_{\ell = 1}^{n}
        {[\underbar{A}]_{i,\ell} [C]_{\ell,1}}
        \\
        = {} &
        \sum_{\ell = 1}^{n}
        {[A]_{i,\ell} [C]_{\ell,1}}
        \\
        = {} & [AC]_{i,1}
        \\
        = {} & [B]_{i,1}
        \\
        = {} &
        % =
        [\underbar{B}]_{i,1}.
    \end{align*}
    若 \(i = k + 1\), 则
    \begin{align*}
        [\underbar{A}C]_{i,1}
        = {} &
        \sum_{\ell = 1}^{n}
        {[\underbar{A}]_{i,\ell} [C]_{\ell,1}}
        \\
        = {} &
        \sum_{\ell = 1}^{n}
        {0\,[C]_{\ell,1}}
        \\
        = {} & 0
        \\
        = {} &
        % =
        [\underbar{B}]_{i,1}.
    \end{align*}
    若 \(i > k + 1\), 则
    \begin{align*}
        [\underbar{A}C]_{i,1}
        = {} &
        \sum_{\ell = 1}^{n}
        {[\underbar{A}]_{i,\ell} [C]_{\ell,1}}
        \\
        = {} &
        \sum_{\ell = 1}^{n}
        {[A]_{i-1,\ell} [C]_{\ell,1}}
        \\
        = {} & [AC]_{i-1,1}
        \\
        = {} & [B]_{i-1,1}
        \\
        = {} &
        % =
        [\underbar{B}]_{i,1}.
    \end{align*}
    由此可见, \(\underbar{A}C = \underbar{B}\).

    (2)
    设 \(\underbar{A}C = \underbar{B}\).
    我们要证 \(AC = B\).

    首先, \(AC\) 与 \(B\)
    的尺寸都是 \(m \times 1\).

    其次, 注意,
    \begin{align*}
        [A]_{i,j}
        = \begin{cases}
              [\underbar{A}]_{i,j},   & i \leq k; \\
              [\underbar{A}]_{i+1,j}, & i > k.
          \end{cases}
    \end{align*}
    类似地,
    \begin{align*}
        [B]_{i,1}
        = \begin{cases}
              [\underbar{B}]_{i,1},   & i \leq k; \\
              [\underbar{B}]_{i+1,1}, & i > k.
          \end{cases}
    \end{align*}

    若 \(i \leq k\), 则
    \begin{align*}
        [AC]_{i,1}
        = {} &
        \sum_{\ell = 1}^{n}
        {[A]_{i,\ell} [C]_{\ell,1}}
        \\
        = {} &
        \sum_{\ell = 1}^{n}
        {[\underbar{A}]_{i,\ell} [C]_{\ell,1}}
        \\
        = {} & [\underbar{A}C]_{i,1}
        \\
        = {} & [\underbar{B}]_{i,1}
        \\
        = {} &
        % =
        [B]_{i,1}.
    \end{align*}
    若 \(i > k\), 则
    \begin{align*}
        [AC]_{i,1}
        = {} &
        \sum_{\ell = 1}^{n}
        {[A]_{i,\ell} [C]_{\ell,1}}
        \\
        = {} &
        \sum_{\ell = 1}^{n}
        {[\underbar{A}]_{i+1,\ell} [C]_{\ell,1}}
        \\
        = {} & [\underbar{A}C]_{i+1,1}
        \\
        = {} & [\underbar{B}]_{i+1,1}
        \\
        = {} &
        % =
        [B]_{i,1}.
    \end{align*}
    由此可见, \(AC = B\).

    (3)
    注意, \(A\)~的子阵都是 \(\underbar{A}\)~的子阵.
    所以, 若 \(A\) 有一个行列式非零的 \(r\)~级子阵,
    则 \(\underbar{A}\) 也有一个行列式非零的 \(r\)~级子阵.
    对 \(\underbar{A}\)~的每个子阵,
    要么 \(\underbar{A}\) 的行~\(k+1\) 被选中,
    要么 \(\underbar{A}\) 的行~\(k+1\) 不被选中.
    所以, 那些 \(\underbar{A}\) 的行~\(k+1\) 不被选中的
    \(r+1\)~级子阵 (若存在)
    是 \(A\)~的 \(r+1\)~级子阵 (若存在),
    故它的行列式为 \(0\).
    而那些 \(\underbar{A}\) 的行~\(k+1\) 被选中的
    \(r+1\)~级子阵 (若存在)
    有一行的元全为 \(0\),
    故它的行列式为 \(0\).
    从而, \(\underbar{A}\) 也没有行列式非零的 \(r+1\)~级子阵.

    (4)
    对 \(\underbar{A}\)~的每个子阵,
    要么 \(\underbar{A}\) 的行~\(k+1\) 被选中,
    要么 \(\underbar{A}\) 的行~\(k+1\) 不被选中.
    设 \(\underbar{A}\) 有一个行列式非零的 \(r\)~级子阵.
    那么, 这个子阵不含元全为 \(0\) 的行.
    特别地, 对此子阵,
    \(\underbar{A}\) 的行~\(k+1\) 不被选中.
    所以, 这也是 \(A\)~的一个 \(r\)~级子阵.
    所以, \(A\) 也有一个行列式非零的 \(r\)~级子阵.
    注意, \(A\)~的子阵都是 \(\underbar{A}\)~的子阵.
    既然 \(\underbar{A}\) 没有行列式非零的 \(r+1\)~级子阵,
    那么 \(A\) 也没有行列式非零的 \(r+1\)~级子阵.
\end{proof}

现在, 我们讨论 \(AX = 0\) 是否有非零解.

\begin{theorem}
    设 \(A\) 是 \(m \times n\)~阵.
    设 \(A\) 有一个行列式非零的 \(r\)~级子阵,
    但没有行列式非零的 \(r+1\)~级子阵.
    设 \(r < n\).
    则 \(AX = 0\) 有非零解.
\end{theorem}

\begin{proof}
    设 \(A\) 是 \(m \times n\)~阵.

    我们无妨设 \(m \geq n\).
    若 \(m < n\), 我们作一个 \(n\)~级阵 \(\underbar{A}\),
    其中,
    \begin{align*}
        [\underbar{A}]_{i,j}
        = \begin{cases}
              [A]_{i,j}, & i \leq m; \\
              0,         & i > m.
          \end{cases}
    \end{align*}
    (通俗地, \(\underbar{A}\)
    是在 \(A\)~的行~\(m\) 下写 \(n-m\)~行 \(0\) 作成的阵.)
    由此, 不难得到:
    (a)
    若 \(n \times 1\)~阵 \(C\) 适合 \(\underbar{A}C = 0\),
    则 \(AC = 0\);
    (b)
    \(\underbar{A}\) 有一个行列式非零的 \(r\)~级子阵,
    但没有行列式非零的 \(r+1\)~级子阵.
    所以, 若我们找到了 \(\underbar{A}X = 0\)
    的非零解,
    则我们也找到了 \(AX = 0\) 的非零解.
    以下, 我们设 \(m \geq n\).

    若 \(A = 0\),
    我们任取一个非零的 \(n \times 1\)~阵 \(S\).
    则 \(AS = 0\).

    以下, 我们假定 \(A \neq 0\);
    也就是, \(A\) 有一个元不是 \(0\)
    (同时, 这也要求 \(n > 1\):
    回想 \(1\)~级阵的行列式是什么).
    从而 \(r \geq 1\).

    我们设 \(A\)~的 \(r\)~级子阵
    \(
    \displaystyle
    A_r = A\binom{i_1,\dots,i_r}{j_1,\dots,j_r}
    \)
    的行列式非零,
    其中, \(i_1 < i_2 < \dots < i_r\),
    且 \(j_1 < j_2 < \dots < j_r\).
    我们从 \(1\), \(2\), \(\dots\), \(m\) 中%
    去除 \(r\)~个整数
    \(i_1\), \(i_2\), \(\dots\), \(i_r\);
    此时, 还剩下 \(m-r\)~个整数;
    我们从中选一个为 \(i_{r+1}\).
    类似地,
    我们再从 \(1\), \(2\), \(\dots\), \(n\) 中%
    去除 \(r\)~个整数
    \(j_1\), \(j_2\), \(\dots\), \(j_r\);
    此时, 还剩下 \(n-r\)~个整数;
    我们从中选一个为 \(j_{r+1}\).
    作 \(r+1\)~级阵
    \(
    \displaystyle
    E =
    A\binom{i_1,\dots,i_r,i_{r+1}}
    {j_1,\dots,j_r,j_{r+1}}
    \).
    我们设 \(i_p\) 在
    \(i_1\), \(\dots\), \(i_r\), \(i_{r+1}\)
    中是第~\(f(i_p)\) 小的数.
    再设 \(j_p\) 在
    \(j_1\), \(\dots\), \(j_r\), \(j_{r+1}\)
    中是第~\(g(j_p)\) 小的数.
    因为 \(E\) 是 \(A\) 的一个 \(r+1\)~级子阵,
    故 \(\det {(E)} = 0\).
    从而
    \(E \operatorname{adj} {(E)} = 0\).
    注意,
    \([\operatorname{adj} {(E)}]_{g(j_{r+1}), f(i_{r+1})}
    = (-1)^{f(i_{r+1})+g(j_{r+1})} \det {(A_r)} \neq 0\).

    接下来, 我们设法由此得到 \(AX = 0\)
    的一个非零解.
    我们作一个 \(n \times 1\)~阵 \(W\), 其中,
    \begin{align*}
        [W]_{\ell,1}
        = \begin{cases}
              [\operatorname{adj} {(E)}]_{g(j_p),f(i_{r+1})},
                 & \text{\(\ell = j_p\),
              \(p = 1\), \(\dots\), \(r\), \(r+1\)}; \\
              0, & \text{别的情形}.
          \end{cases}
    \end{align*}
    (通俗地, 我们在适当的位置写 \(0\),
    变 \(\operatorname{adj} {(E)}\)
    的列~\(f(i_{r+1})\)
    为一个 \(n \times 1\)~阵 \(W\).)
    因为 \([W]_{j_{r+1},1} \neq 0\),
    故 \(W \neq 0\).
    我们证明 \(AW = 0\).

    取不超过 \(m\) 的正整数 \(q\).
    则
    \begin{align*}
        [AW]_{q,1}
        = {} &
        \sum_{\ell = 1}^{n}
        {[A]_{q,\ell} [W]_{\ell,1}}
        \\
        = {} &
        \sum_{\substack{
        1 \leq \ell \leq n \\
        \ell = j_p         \\
                1 \leq p \leq r+1
            }}
        {[A]_{q,\ell} [W]_{\ell,1}}
        +
        \sum_{\substack{
        1 \leq \ell \leq n \\
                \text{别的情形}
            }}
        {[A]_{q,\ell} [W]_{\ell,1}}
        \\
        = {} &
        \sum_{\substack{
        1 \leq \ell \leq n \\
        \ell = j_p         \\
                1 \leq p \leq r+1
            }}
        {[A]_{q,j_p}
                [\operatorname{adj} {(E)}]_{g(j_p),f(i_{r+1})}}
        +
        \sum_{\substack{
        1 \leq \ell \leq n \\
                \text{别的情形}
            }}
        {[A]_{q,\ell}\, 0}
        \\
        = {} &
        \sum_{p=1}^{r+1}
        {[A]_{q,j_p}
            [\operatorname{adj} {(E)}]_{g(j_p),f(i_{r+1})}}.
    \end{align*}
    若 \(q\) 等于某个 \(i_t\), 则
    \begin{align*}
        [AW]_{i_t,1}
        = {} &
        \sum_{p=1}^{r+1}
        {[A]_{i_t,j_p}
            [\operatorname{adj} {(E)}]_{g(j_p),f(i_{r+1})}}
        \\
        = {} &
        \sum_{p=1}^{r+1}
        {[E]_{f(i_t),g(j_p)}
            [\operatorname{adj} {(E)}]_{g(j_p),f(i_{r+1})}}
        \\
        = {} &
        [E \operatorname{adj} {(E)}]_{f(i_t),f(i_{r+1})}
        \\
        = {} &
        [0]_{f(i_t),f(i_{r+1})}
        \\
        = {} & 0.
    \end{align*}
    若 \(q\) 不等于
    \(i_1\), \(\dots\), \(i_r\), \(i_{r+1}\)
    中的任何一个,
    则
    \begin{align*}
        [AW]_{q,1}
        = {} &
        \sum_{p=1}^{r+1}
        {[A]_{q,j_p}
            [\operatorname{adj} {(E)}]_{g(j_p),f(i_{r+1})}}
        \\
        = {} &
        \sum_{p=1}^{r+1}
        {[A]_{q,j_p}
        (-1)^{f(i_{r+1})+g(j_p)}
        \det {(E(f(i_{r+1})|g(j_p)))}}.
    \end{align*}
    作一个 \(r+1\)~级阵 \(C_q\), 其中,
    \begin{align*}
        [C_q]_{h,g(j_p)}
        = \begin{cases}
              [E]_{h,g(j_p)},
               & h \neq f(i_{r+1}); \\
              [A]_{q,j_p},
               & h = f(i_{r+1}).
          \end{cases}
    \end{align*}
    于是, \(C_q(f(i_{r+1})|g(j_p))
    = E(f(i_{r+1})|g(j_p))\).
    从而
    \begin{align*}
        [AW]_{q,1}
        = {} &
        \sum_{p=1}^{r+1}
        {[A]_{q,j_p}
        (-1)^{f(i_{r+1})+g(j_p)}
        \det {(E(f(i_{r+1})|g(j_p)))}}
        \\
        = {} &
        \sum_{p=1}^{r+1}
        {[C_q]_{f(i_{r+1}),g(j_p)}
        (-1)^{f(i_{r+1})+g(j_p)}
        \det {(C_q (f(i_{r+1})|g(j_p)))}}
        \\
        = {} &
        \det {(C_q)}.
    \end{align*}

    再作一个 \(r+1\)~级阵
    \(
    \displaystyle
    D_q =
    A\binom{i_1,\dots,i_r,q}
    {j_1,\dots,j_r,j_{r+1}}
    \).
    不难看出, 适当地交换 \(C_q\) 的行的次序,
    即可变 \(C_q\) 为 \(D_q\).
    根据反称性,
    \(\det {(C_q)} = \pm \det {(D_q)}\).
    (具体地, 设 \(q\) 在
    \(i_1\), \(\dots\), \(i_r\), \(q\)
    中是第~\(u\) 小的数.
    那么, 当 \(f(i_{r+1}) = u\) 时, \(C_q = D_q\).
    当 \(f(i_{r+1}) < u\) 时,
    交换行~\(f(i_{r+1})\) 与 \(f(i_{r+1})+1\),
    再交换行~\(f(i_{r+1})+1\) 与 \(f(i_{r+1})+2\),
    ……
    且再交换行~\(u-1\) 与 \(u\);
    作 \(u - f(i_{r+1})\)~次相邻行的交换,
    即可变 \(C_q\) 为 \(D_q\).
    当 \(f(i_{r+1}) > u\) 时,
    交换行~\(f(i_{r+1})\) 与 \(f(i_{r+1})-1\),
    再交换行~\(f(i_{r+1})-1\) 与 \(f(i_{r+1})-2\),
    ……
    且再交换行~\(u+1\) 与 \(u\);
    作 \(f(i_{r+1}) - u\)~次相邻行的交换,
    即可变 \(C_q\) 为 \(D_q\).)
    因为 \(D_q\) 是 \(A\) 的一个 \(r+1\)~级子阵,
    故它的行列式为 \(0\).
    从而
    \begin{align*}
        [AW]_{q,1}
        = \det {(C_q)}
        = \pm \det {(D_q)}
        = 0.
    \end{align*}

    综上, 我们找到了一个 \(n \times 1\)~阵 \(W\)
    使 \(AW = 0\), 且 \(W \neq 0\).
\end{proof}

\begin{theorem}
    设 \(A\) 是 \(m \times n\)~阵.
    设 \(A\) 有一个行列式非零的 \(n\)~级子阵
    (此时, \(A\) 当然没有行列式非零的 \(n+1\)~级子阵).
    则 \(AX = 0\) 只有零解.
\end{theorem}

\begin{proof}
    设 \(n \times 1\)~阵 \(C\) 适合 \(AC = 0\).
    设 \(A\)~的 \(n\)~级子阵
    \begin{align*}
        K = A\binom{i_1,i_2,\dots,i_n}{1,2,\dots,n}
    \end{align*}
    的行列式非零,
    其中,
    \(1 \leq i_1 < i_2 < \dots < i_n \leq m\).
    因为 \(C\) 适合
    \begin{align*}
        [A]_{1,1} [C]_{1,1} + [A]_{1,2} [C]_{2,1} + \dots + [A]_{1,n} [C]_{n,1} & = 0,   \\
        [A]_{2,1} [C]_{1,1} + [A]_{2,2} [C]_{2,1} + \dots + [A]_{2,n} [C]_{n,1} & = 0,   \\
                                                                                & \dots, \\
        [A]_{m,1} [C]_{1,1} + [A]_{m,2} [C]_{2,1} + \dots + [A]_{m,n} [C]_{n,1} & = 0,
    \end{align*}
    故 \(C\) 当然也适合
    \begin{align*}
        [A]_{i_1,1} [C]_{1,1} + [A]_{i_1,2} [C]_{2,1} + \dots + [A]_{i_1,n} [C]_{n,1} & = 0,   \\
        [A]_{i_2,1} [C]_{1,1} + [A]_{i_2,2} [C]_{2,1} + \dots + [A]_{i_2,n} [C]_{n,1} & = 0,   \\
                                                                                      & \dots, \\
        [A]_{i_n,1} [C]_{1,1} + [A]_{i_n,2} [C]_{2,1} + \dots + [A]_{i_n,n} [C]_{n,1} & = 0,
    \end{align*}
    即 \(KC = 0\).
    因为 \(\det {(K)} \neq 0\),
    故, 由 Cramer 公式,
    有 \(C = 0\).
\end{proof}

现在, 我们作一个小结.

\begin{theorem}
    设 \(A\) 是 \(m \times n\)~阵.
    设 \(A\) 有一个行列式非零的 \(r\)~级子阵,
    但没有行列式非零的 \(r+1\)~级子阵.

    (1)
    若 \(r < n\),
    则 \(AX = 0\) 有非零解;

    (2)
    若 \(r = n\),
    则 \(AX = 0\) 只有零解.
\end{theorem}

\begin{theorem}
    设 \(A\) 是 \(m \times n\)~阵.
    设 \(B\) 是 \(m \times 1\)~阵.
    设 \(n \times 1\)~阵 \(C\) 适合 \(AC = B\);
    也就是, \(AX = B\) 有解.
    设 \(A\) 有一个行列式非零的 \(r\)~级子阵,
    但没有行列式非零的 \(r+1\)~级子阵.

    (1)
    若 \(r < n\),
    则 \(AX = B\) 的解不唯一;

    (2)
    若 \(r = n\),
    则 \(AX = B\) 的解唯一.
\end{theorem}

\KunAsteriskoEnEnhavtabelo
\section{\texorpdfstring{由 \(m\)~个 \(n\)~元
      \({\leq} 1\)~次方程作成的方程组 (2)}%
  {由 m 个 n 元 ≤1 次方程作成的方程组 (2)}}
\SenAsteriskoEnEnhavtabelo

\maldevigalegajxo

前面, 我们讨论了
``\(AX = B\) 有解时, 解是否唯一''
问题.
此事的回答跟 \(A\)~的子阵的行列式有关.
设 \(A\) 有 \(n\)~列.
当 \(A\) 有一个行列式非零的 \(n\)~级子阵时,
此方程组的解唯一;
当 \(A\) 没有行列式非零的 \(n\)~级子阵时,
此方程组的解不唯一.
本节, 我们讨论线性方程组何时有解.
一个好的想法是,
我们先讨论 \(AX = B\) 有解时,
\(A\), \(B\) 应适合什么条件;
然后, 反过来,
我们再讨论适合这些条件的 \(A\), \(B\)
是否能使 \(AX = B\) 有解.

% restatable
\begin{restatable}{theorem}{TheoremNecessityForConsistency}
    设 \(A\) 是 \(m \times n\)~阵.
    设 \(B\) 是 \(m \times 1\)~阵.
    设存在 \(n \times 1\)~阵 \(C\)
    适合 \(AC = B\).
    作一个 \(m \times (n+1)\)~阵 \(G\),
    其中,
    \begin{align*}
        [G]_{i,j}
        = \begin{cases}
              [A]_{i,j}, & j \leq n;  \\
              [B]_{i,1}, & j = n + 1.
          \end{cases}
    \end{align*}
    (通俗地, 在 \(A\)~的最后一列的右侧加入一列 \(B\),
    得到尺寸较大的阵 \(G\).)
    则存在一个非负整数 \(r\), 使
    \(A\) 有一个行列式非零的 \(r\)~级子阵
    (从而 \(G\) 也有一个行列式非零的 \(r\)~级子阵),
    但 \(G\) 没有行列式非零的 \(r+1\)~级子阵
    (从而 \(A\) 也没有行列式非零的 \(r+1\)~级子阵).
\end{restatable}

\begin{proof}
    我们知道, 存在一个非负整数 \(r\),
    使 \(A\) 有一个行列式非零的 \(r\)~级子阵,
    但 \(A\) 没有行列式非零的 \(r+1\)~级子阵.
    我们证明,
    \(G\) 也没有行列式非零的 \(r+1\)~级子阵.

    若 \(G\) 根本没有 \(r+1\)~级子阵,
    则 \(G\) 当然也没有行列式非零的 \(r+1\)~级子阵.

    若 \(G\) 有 \(r+1\)~级子阵,
    那么,
    对 \(G\)~的每个 \(r+1\)~级子阵,
    要么 \(G\) 的列~\(n+1\) 被选中,
    要么 \(G\) 的列~\(n+1\) 不被选中.

    若 \(G\) 的列~\(n+1\) 不被选中,
    那么, 这个子阵就是 \(A\)~的 \(r+1\)~级子阵,
    故它的行列式为 \(0\).

    若 \(G\) 的列~\(n+1\) 被选中,
    我们说明,
    这个子阵的行列式是
    \(A\)~的一些 \(r+1\)~级子阵的行列式的倍的和,
    故它的行列式仍为 \(0\).
    因为 \(AC = B\), 故,
    对任何不超过 \(m\) 的正整数 \(i\),
    \begin{align*}
        [B]_{i,1}
            = [A]_{i,1} [C]_{1,1} + [A]_{i,2} [C]_{2,1}
        + \dots + [A]_{i,n} [C]_{n,1},
    \end{align*}
    即
    \begin{align*}
        [G]_{i,n+1}
        = [G]_{i,1} [C]_{1,1} + [G]_{i,2} [C]_{2,1}
        + \dots + [G]_{i,n} [C]_{n,1}.
    \end{align*}
    设这个子阵为
    \begin{align*}
        D = G\binom{i_1,\dots,i_r,i_{r+1}}{j_1,\dots,j_r,n+1},
    \end{align*}
    其中,
    \(1 \leq i_1 < \dots < i_r < i_{r+1} \leq m\),
    且 \(1 \leq j_1 < \dots < j_r \leq n\).
    记 \((r+1) \times n\)~阵
    \begin{align*}
        A\binom{i_1,\dots,i_r,i_{r+1}}{1,2,\dots,n}
    \end{align*}
    的列~\(1\), \(2\), \(\dots\), \(n\) 为
    \(h_1\), \(h_2\), \(\dots\), \(h_n\).
    于是, \(D\)~的列~\(1\), \(2\), \(\dots\), \(r\) 就是
    \(h_{j_1}\), \(h_{j_2}\), \(\dots\), \(h_{j_r}\),
    而 \(D\)~的列~\(r+1\) 是
    \begin{align*}
        \sum_{1 \leq t \leq n}
        {[C]_{t,1} h_t}.
    \end{align*}
    从而, 利用多线性,
    \begin{align*}
        \det {(D)}
        = {} &
        \det {
            \left[
                h_{j_1}, \dots, h_{j_r},
                \sum_{1 \leq t \leq n}
                {[C]_{t,1} h_t}
                \right]
        }
        \\
        = {} &
        \sum_{1 \leq t \leq n}
        {
            [C]_{t,1} \det {[h_{j_1}, \dots, h_{j_r}, h_t]}
        }.
    \end{align*}
    若 \(t\) 等于 \(j_1\), \(j_2\), \(\dots\), \(j_r\) 的一个,
    则 \([h_{j_1}, \dots, h_{j_r}, h_t]\) 有二列相同.
    由交错性, \(\det {[h_{j_1}, \dots, h_{j_r}, h_t]} = 0\).
    若 \(t\) 不等于 \(j_1\), \(j_2\), \(\dots\), \(j_r\) 的任何一个,
    我们适当地交换 \([h_{j_1}, \dots, h_{j_r}, h_t]\) 的列的次序,
    以变它为 \(A\) 的一个 \(r+1\)~级子阵.
    由反称性, 有
    \begin{align*}
        \det {[h_{j_1}, \dots, h_{j_r}, h_t]}
        = \pm
        \det {\left(
            A\binom{i_1,\dots,i_r,i_{r+1}}{j_1,\dots,j_r,t}
            \right)}
        = 0.
    \end{align*}
    于是, \(\det {[h_{j_1}, \dots, h_{j_r}, h_t]}\) 总是 \(0\).
    故 \(\det {(D)} = 0\).
\end{proof}

\KunAsteriskoEnEnhavtabelo
\section{\texorpdfstring{由 \(m\)~个 \(n\)~元
      \({\leq} 1\)~次方程作成的方程组 (3)}%
  {由 m 个 n 元 ≤1 次方程作成的方程组 (3)}}
\SenAsteriskoEnEnhavtabelo

\maldevigalegajxo

前面, 我们讨论了
\(AX = B\) 有解时,
\(A\), \(B\) 应适合的条件:

\TheoremNecessityForConsistency*

那么, 反过来,
若 \(A\) 有一个行列式非零的 \(r\)~级子阵,
但 \(G\) 没有行列式非零的 \(r+1\)~级子阵,
则 \(AX = B\) 是否有解?
此事的回答是 ``是''.
不过, 为了论证此事,
我们要作准备.

\begin{theorem}
    设 \(A\) 是 \(n\)~级阵.
    设 \(p\), \(q\) 是二个不超过 \(n\)~的正整数,
    且 \(p \neq q\).
    设 \(x\) 是一个数.
    作 \(n\)~级阵 \(L\), 其中,
    \begin{align*}
        [L]_{i,j}
        = \begin{cases}
              [A]_{i,j},              & j \neq q; \\
              [A]_{i,q} + x[A]_{i,p}, & j = q.
          \end{cases}
    \end{align*}
    (通俗地,
    加 \(A\)~的列~\(p\) 的 \(x\)~倍于列~\(q\),
    且不改变别的列,
    得阵~\(L\).)
    则 \(\det {(L)} = \det {(A)}\).

    类似地,
    若加 \(A\)~的行~\(p\) 的 \(x\)~倍于行~\(q\),
    且不改变别的行,
    得阵~\(K\),
    则 \(\det {(K)} = \det {(A)}\).
\end{theorem}

\begin{proof}
    我证明关于列的事;
    我们可用类似的方法证明关于行的事
    (或者, 利用行列式与转置的关系).

    作一个 \(n\)~级阵 \(S\),
    使 \(S\) 的列~\(j\) 等于 \(A\) 的列~\(j\)
    (若 \(j \neq q\)),
    且 \(S\) 的列~\(q\) 等于 \(A\) 的列~\(p\).
    则 \(S\) 的列~\(p\) 等于 \(S\) 的列~\(q\),
    且 \(L\) 的列~\(j\), \(A\) 的列~\(j\),
    与 \(S\) 的列~\(j\) 相等
    (若 \(j \neq q\)),
    但 \([L]_{i,q} = 1 [A]_{i,q} + x [S]_{i,q}\).
    由多线性与交错性,
    \begin{equation*}
        \det {(L)}
        % = {} &
        =
        1 \det {(A)} + x \det {(S)}
        % \\
        % = {} &
        =
        \det {(A)} + x\, 0
        % \\
        % = {} &
        =
        \det {(A)}.
        \qedhere
    \end{equation*}
    % The old proof that is unnecessarily lengthy.
    % 设 \(A = [a_1, a_2, \dots, a_n]\).

    % 先设 \(p < q\).
    % 为方便, 我们写
    % \begin{align*}
    %     f(u, v)
    %     = \det {[a_1, \dots, a_{p-1}, u, a_{p+1}, \dots,
    %                 a_{q-1}, v, a_{q+1}, \dots, a_n]}.
    % \end{align*}
    % 于是, \(f(a_p, a_q)\) 就是 \(\det {(A)}\),
    % 而 \(f(a_p, a_q + xa_p)\) 就是 \(\det {(L)}\).
    % 利用多线性与交错性,
    % \begin{align*}
    %     \det {(L)}
    %     = {} & f(a_p, a_q + xa_p)         \\
    %     = {} & f(a_p, a_q) + xf(a_p, a_p) \\
    %     = {} & f(a_p, a_q) + x0           \\
    %     = {} & f(a_p, a_q)                \\
    %     = {} & \det {(A)}.
    % \end{align*}

    % 再设 \(p > q\).
    % 为方便, 我们写
    % \begin{align*}
    %     g(u, v)
    %     = \det {[a_1, \dots, a_{q-1}, u, a_{q+1}, \dots,
    %                 a_{p-1}, v, a_{p+1}, \dots, a_n]}.
    % \end{align*}
    % 于是, \(g(a_q, a_p)\) 就是 \(\det {(A)}\),
    % 而 \(g(a_q + xa_p, a_p)\) 就是 \(\det {(L)}\).
    % 利用多线性与交错性,
    % \begin{align*}
    %     \det {(L)}
    %     = {} & g(a_q + xa_p, a_p)         \\
    %     = {} & g(a_q, a_p) + xg(a_p, a_p) \\
    %     = {} & g(a_q, a_p) + x0           \\
    %     = {} & g(a_q, a_p)                \\
    %     = {} & \det {(A)}.
    %     \qedhere
    % \end{align*}
\end{proof}

\begin{theorem}
    设 \(A\) 是 \(m \times n\)~阵,
    且 \(A \neq 0\).
    设 \(A\) 有一个行列式非零的 \(r\)~级子阵
    \begin{align*}
        A_r = A\binom{i_1,\dots,i_r}{j_1,\dots,j_r}
    \end{align*}
    (其中,
    \(1 \leq i_1 < \dots < i_r \leq m\),
    且 \(1 \leq j_1 < \dots < j_r \leq n\)),
    但 \(A\) 没有行列式非零的 \(r+1\)~级子阵.
    那么, 对任何不超过 \(m\) 的正整数 \(p\),
    存在 \(r\)~个数
    \(d_{p,1}\), \(d_{p,2}\), \(\dots\), \(d_{p,r}\),
    使对任何不超过 \(n\) 的正整数 \(q\),
    \begin{align*}
        [A]_{p,q}
        % = {} &
            =
            [A]_{i_1,q} d_{p,1}
        + [A]_{i_2,q} d_{p,2}
        + \dots
        + [A]_{i_r,q} d_{p,r}
        % \\
        % = {} &
        =
        \sum_{s = 1}^{r} {[A]_{i_s,q} d_{p,s}}.
    \end{align*}
\end{theorem}

我们也可如此说前面的结论:

\begin{theorem}
    设 \(A\) 是 \(m \times n\)~阵,
    且 \(A \neq 0\).
    设 \(A\) 有一个行列式非零的 \(r\)~级子阵
    \begin{align*}
        A_r = A\binom{i_1,\dots,i_r}{j_1,\dots,j_r}
    \end{align*}
    (其中,
    \(1 \leq i_1 < \dots < i_r \leq m\),
    且 \(1 \leq j_1 < \dots < j_r \leq n\)),
    但 \(A\) 没有行列式非零的 \(r+1\)~级子阵.
    设 \(A\)~的%
    行~\(1\), \(2\), \(\dots\), \(m\)
    为 \(a_1\), \(a_2\), \(\dots\), \(a_m\).
    那么, 对任何不超过 \(m\) 的正整数 \(p\),
    存在 \(r\)~个数
    \(d_{p,1}\), \(d_{p,2}\), \(\dots\), \(d_{p,r}\),
    使
    \begin{align*}
        a_p
        % = {} &
        =
        d_{p,1} a_{i_1}
        + d_{p,2} a_{i_2}
        + \dots
        + d_{p,r} a_{i_r}
        % \\
        % = {} &
        =
        \sum_{s = 1}^{r} {d_{p,s} a_{i_s}}.
    \end{align*}
\end{theorem}

通俗地, 这个定理说,
任给一个非零阵 \(A\),
我们总能找出它的某 \(r\)~行
\(a_{i_1}\), \(a_{i_2}\), \(\dots\), \(a_{i_r}\),
使 \(A\)~的每行都可被写为%
这 \(r\)~行的数乘的和,
且由这 \(r\)~行作成的子阵有一个%
行列式非零的 \(r\)~级子阵.

\begin{example}
    设
    \begin{align*}
        A =
        \begin{bmatrix}
            1 & 0 & 0 & 0 & 0 & -1 \\
            0 & 1 & 0 & 0 & 0 & 1  \\
            0 & 0 & 1 & 0 & 0 & 0  \\
            0 & 1 & 1 & 0 & 0 & 1  \\
        \end{bmatrix}.
    \end{align*}
    不难看出, \(A\) 有一个 \(3\)~级子阵, 其的行列式非零:
    \begin{align*}
        \det {\left(
            A\binom{1,2,3}{1,2,3}
            \right)} = 1.
    \end{align*}
    不过, \(A\) 没有行列式非零的 \(4\)~级子阵.
    我们考虑 \(A\) 的列~\(4\), \(5\).
    任取 \(A\)~的一个 \(4\)~级子阵.
    若 \(A\) 的列~\(4\) 或列~\(5\) 被选中,
    那么它的行列式显然为 \(0\).
    若 \(A\) 的列~\(4\) 与列~\(5\) 都不被选中,
    则这个子阵是
    \begin{align*}
        A\binom{1,2,3,4}{1,2,3,6}
        = \begin{bmatrix}
              1 & 0 & 0 & -1 \\
              0 & 1 & 0 & 1  \\
              0 & 0 & 1 & 0  \\
              0 & 1 & 1 & 1  \\
          \end{bmatrix}.
    \end{align*}
    不难算出, 它的行列式为 \(0\).

    我们说, \(A\) 的每行可被写为
    \(A\) 的前 \(3\)~行的数乘的和.
    设 \(A\)~的%
    行~\(1\), \(2\), \(3\), \(4\)
    为 \(a_1\), \(a_2\), \(a_3\), \(a_4\).
    则
    \begin{align*}
        a_1 & = 1a_1 + 0a_2 + 0a_3, \\
        a_2 & = 0a_1 + 1a_2 + 0a_3, \\
        a_3 & = 0a_1 + 0a_2 + 1a_3, \\
        a_4 & = 0a_1 + 1a_2 + 1a_3.
    \end{align*}
    并且, 由这 \(3\)~行作成的子阵有一个%
    行列式非零的 \(3\)~级子阵:
    \(3\)~级单位阵就是一个.

    当然,
    \(A\) 的每行可被写为
    \(A\) 的前 \(4\)~行的数乘的和:
    \begin{align*}
        a_1 & = 1a_1 + 0a_2 + 0a_3 + 0a_4, \\
        a_2 & = 0a_1 + 1a_2 + 0a_3 + 0a_4, \\
        a_3 & = 0a_1 + 0a_2 + 1a_3 + 0a_4, \\
        a_4 & = 0a_1 + 0a_2 + 0a_3 + 1a_4.
    \end{align*}
    可是, 由这 \(4\)~行作成的子阵没有%
    行列式非零的 \(4\)~级子阵.
\end{example}

\begin{proof}
    任取不超过 \(m\) 的正整数 \(p\).

    若 \(p\) 等于某个 \(i_s\)
    (\(s = 1\), \(2\), \(\dots\), \(r\)),
    我们取
    \begin{align*}
        d_{p,v}
        = \begin{cases}
              1, & v = s;    \\
              0, & v \neq s.
          \end{cases}
    \end{align*}
    于是, 对任何不超过 \(n\) 的正整数 \(q\),
    \begin{align*}
        [A]_{i_1,q} d_{p,1}
        + [A]_{i_2,q} d_{p,2}
        + \dots
        + [A]_{i_r,q} d_{p,r}
            = [A]_{i_s,q} 1
            = [A]_{p,q}.
    \end{align*}

    下设 \(p\) 不等于任何 \(i_s\).

    考虑%
    由 \(r\)~个 \(r\)~元 \({\leq} 1\)~次方程作成的方程组
    \begin{equation*}
        \left\{
        \begin{aligned}
            [A]_{i_1,j_1} x_1 + [A]_{i_2,j_1} x_2 + \dots + [A]_{i_r,j_1} x_r & = [A]_{p,j_1}, \\
            [A]_{i_1,j_2} x_1 + [A]_{i_2,j_2} x_2 + \dots + [A]_{i_r,j_2} x_r & = [A]_{p,j_2}, \\
                                                                              & \dots,         \\
            [A]_{i_1,j_r} x_1 + [A]_{i_2,j_r} x_2 + \dots + [A]_{i_r,j_r} x_r & = [A]_{p,j_r},
        \end{aligned}
        \right.
    \end{equation*}
    即
    \begin{align*}
        A_r^{\mathrm{T}}
        \begin{bmatrix}
            x_1    \\
            x_2    \\
            \vdots \\
            x_r    \\
        \end{bmatrix}
        =
        \begin{bmatrix}
            [A]_{p,j_1} \\
            [A]_{p,j_2} \\
            \vdots      \\
            [A]_{p,j_r} \\
        \end{bmatrix}.
    \end{align*}
    因为
    \(\det {(A_r^{\mathrm{T}})}
    = \det {(A_r)} \neq 0\),
    故, 由 Cramer 公式,
    存在 \(r\)~个数
    \(d_{p,1}\), \(d_{p,2}\), \(\dots\), \(d_{p,r}\),
    使
    \begin{align*}
        [A]_{i_1,j_1} d_{p,1} + [A]_{i_2,j_1} d_{p,2} + \dots + [A]_{i_r,j_1} d_{p,r} & = [A]_{p,j_1}, \\
        [A]_{i_1,j_2} d_{p,1} + [A]_{i_2,j_2} d_{p,2} + \dots + [A]_{i_r,j_2} d_{p,r} & = [A]_{p,j_2}, \\
                                                                                      & \dots,         \\
        [A]_{i_1,j_r} d_{p,1} + [A]_{i_2,j_r} d_{p,2} + \dots + [A]_{i_r,j_r} d_{p,r} & = [A]_{p,j_r}.
    \end{align*}
    我们由此证明, 对任何不超过 \(n\) 的正整数 \(q\),
    \begin{align*}
        [A]_{p,q}
            =
            [A]_{i_1,q} d_{p,1}
        + [A]_{i_2,q} d_{p,2}
        + \dots
        + [A]_{i_r,q} d_{p,r}.
    \end{align*}

    若 \(q\) 等于某个 \(j_\ell\), 显然.
    下设 \(q\) 不等于任何 \(j_\ell\).

    考虑 \(A\) 的 \(r+1\)~级子阵
    \begin{align*}
        K =
        A\binom{i_1,\dots,i_r,p}{j_1,\dots,j_r,q}.
    \end{align*}
    我们用重要思想, 算二次,
    证我们想要的等式.

    一方面, 我们知道, 既然
    \(A\) 没有行列式非零的 \(r+1\)~级子阵,
    故 \(\det {(K)} = 0\).

    另一方面, 我们也可适当地作辅助阵,
    其的行列式等于 \(K\) 的行列式.
    从而, 计算这些辅助阵的行列式,
    也就相当于计算 \(K\) 的行列式.
    为方便, 我们记
    \(i_{r+1} = p\),
    且 \(j_{r+1} = q\).
    设 \(i_s\) 在
    \(i_1\), \(\dots\), \(i_r\), \(i_{r+1}\)
    中是第~\(f(i_s)\) 小的数;
    设 \(j_\ell\) 在
    \(j_1\), \(\dots\), \(j_r\), \(j_{r+1}\)
    中是第~\(g(j_\ell)\) 小的数.
    加 \(K\) 的行~\(f(i_1)\) 的 \(-d_{p,1}\)~倍%
    于行~\(f(p)\),
    得阵~\(K_1\).
    则 \(\det {(K_1)} = \det {(K)}\).
    加 \(K_1\) 的行~\(f(i_2)\) 的 \(-d_{p,2}\)~倍%
    于行~\(f(p)\),
    得阵~\(K_2\).
    则 \(\det {(K_2)} = \det {(K_1)} = \det {(K)}\).
    ……
    加 \(K_{r-1}\) 的行~\(f(i_r)\) 的 \(-d_{p,r}\)~倍%
    于行~\(f(p)\),
    得阵~\(K_r\).
    则 \(\det {(K_r)} = \det {(K_{r-1})} = \det {(K)}\).
    注意, \(K_r\)~的行
    \(f(i_1)\), \(f(i_2)\), \(\dots\), \(f(i_r)\)
    分别跟 \(K\)~的行
    \(f(i_1)\), \(f(i_2)\), \(\dots\), \(f(i_r)\)
    相等;
    不过,
    \begin{align*}
        [K_r]_{f(p),g(v)}
        = [A]_{p,v} -
        \sum_{s = 1}^{r}
        {[A]_{i_s,v} d_{p,s}},
    \end{align*}
    其中, \(v = j_1\), \(\dots\), \(j_r\), \(q\).
    所以, 当
    % \(t \neq r+1\)
    \(g(v) \neq g(q)\)
    时,
    % \([K_r]_{f(p),g(j_t)} = 0\).
    \([K_r]_{f(p),g(v)} = 0\).
    我们按行~\(f(p)\) 展开 \(K_r\) 的行列式,
    有
    \begin{align*}
        \det {(K_r)}
        = {} &
        (-1)^{f(p)+g(q)} [K_r]_{f(p),g(q)}
        \det {(K_r (f(p)|g(q)))}
        \\
        = {} &
        (-1)^{f(p)+g(q)} \det {(A_r)}
        \left(
        [A]_{p,q} -
        \sum_{s = 1}^{r}
        {[A]_{i_s,q} d_{p,s}}
        \right).
    \end{align*}

    回想, \(0 = \det {(K)}\),
    且 \(\det {(K)} = \det {(K_r)}\).
    比较二次计算的结果, 我们应有
    \begin{align*}
        0 =
        (-1)^{f(p)+g(q)} \det {(A_r)}
        \left(
        [A]_{p,q} -
        \sum_{s = 1}^{r}
        {[A]_{i_s,q} d_{p,s}}
        \right).
    \end{align*}
    注意,
    \((-1)^{f(p)+g(q)} \det {(A_r)} \neq 0\),
    故
    \begin{equation*}
        [A]_{p,q} -
        \sum_{s = 1}^{r}
        {[A]_{i_s,q} d_{p,s}} = 0.
        \qedhere
    \end{equation*}
\end{proof}

现在, 我们可以证明本节的重要结论了.

\begin{theorem}
    设 \(A\) 是 \(m \times n\)~阵.
    设 \(B\) 是 \(m \times 1\)~阵.
    作一个 \(m \times (n+1)\)~阵 \(G\),
    其中,
    \begin{align*}
        [G]_{i,j}
        = \begin{cases}
              [A]_{i,j}, & j \leq n;  \\
              [B]_{i,1}, & j = n + 1.
          \end{cases}
    \end{align*}
    (通俗地, 在 \(A\)~的最后一列的右侧加入一列 \(B\),
    得到尺寸较大的阵 \(G\).)
    设
    \(A\) 有一个行列式非零的 \(r\)~级子阵
    (从而 \(G\) 也有一个行列式非零的 \(r\)~级子阵),
    但 \(G\) 没有行列式非零的 \(r+1\)~级子阵
    (从而 \(A\) 也没有行列式非零的 \(r+1\)~级子阵).
    则存在一个 \(n \times 1\)~阵 \(C\),
    使 \(AC = B\).
\end{theorem}

\begin{proof}
    先设 \(A = 0\).
    那么, \(A\) 有一个行列式非零的 ``\(0\)~级子阵'',
    但 \(G\) 没有行列式非零的 \(1\)~级子阵.
    所以 \(G = 0\).
    所以 \(B = 0\).
    那么, 任何 \(n \times 1\)~阵 \(C\)
    都适合 \(AC = B\).

    下设 \(A \neq 0\).
    设 \(A\) 的 \(r\)~级子阵
    \begin{align*}
        A_r = A\binom{i_1,\dots,i_r}{j_1,\dots,j_r}
    \end{align*}
    (其中,
    \(1 \leq i_1 < \dots < i_r \leq m\),
    且 \(1 \leq j_1 < \dots < j_r \leq n\))
    的行列式非零.
    \(A_r\) 当然也是 \(G\)~的子阵, 且
    \begin{align*}
        A\binom{i_1,\dots,i_r}{j_1,\dots,j_r}
        =
        G\binom{i_1,\dots,i_r}{j_1,\dots,j_r}.
    \end{align*}
    所以, 对任何不超过 \(m\) 的正整数 \(p\),
    存在 \(r\)~个数
    \(d_{p,1}\), \(d_{p,2}\), \(\dots\), \(d_{p,r}\),
    使对任何不超过 \(n+1\) 的正整数 \(q\),
    \begin{align*}
        [G]_{p,q}
        % = {} &
            =
            [G]_{i_1,q} d_{p,1}
        + [G]_{i_2,q} d_{p,2}
        + \dots
        + [G]_{i_r,q} d_{p,r}
        % \\
        % = {} &
        =
        \sum_{s = 1}^{r} {[G]_{i_s,q} d_{p,s}}.
    \end{align*}

    考虑%
    由 \(r\)~个 \(n\)~元 \({\leq} 1\)~次方程作成的方程组
    \begin{equation*}
        \left\{
        \begin{aligned}
            [A]_{i_1,1} x_1 + [A]_{i_1,2} x_2 + \dots + [A]_{i_1,n} x_n & = [B]_{i_1,1}, \\
            [A]_{i_2,1} x_1 + [A]_{i_2,2} x_2 + \dots + [A]_{i_2,n} x_n & = [B]_{i_2,1}, \\
                                                                        & \dots,         \\
            [A]_{i_r,1} x_1 + [A]_{i_r,2} x_2 + \dots + [A]_{i_r,n} x_n & = [B]_{i_r,1}.
        \end{aligned}
        \right.
    \end{equation*}
    我们说, 这个方程组有解.

    若 \(r = n\), 则这是一个%
    由 \(n\)~个 \(n\)~元 \({\leq} 1\)~次方程作成的方程组.
    因为 \(\det {(A_r)} \neq 0\),
    故, 由 Cramer 公式,
    此方程组有一个 (唯一的) 解.

    若 \(r < n\),
    从 \(1\), \(2\), \(\dots\), \(n\)
    去除 \(j_1\), \(j_2\), \(\dots\), \(j_r\)
    后, 还剩 \(n - r\)~个数.
    我们从小到大地叫这 \(n - r\)~个数为
    \(j_{r+1}\), \(\dots\), \(j_n\).
    我们可改写此方程组为
    \begin{align*}
        \sum_{\ell = 1}^{r}
        {[A]_{p,j_\ell} x_{j_\ell}}
            = [B]_{p,1}
        - \sum_{r < \ell \leq n}
        {[A]_{p,j_\ell} x_{j_\ell}},
    \end{align*}
    其中, \(p = i_1\), \(i_2\), \(\dots\), \(i_r\),
    下同.

    设
    \(f_{j_{r+1}}\), \(\dots\), \(f_{j_r}\)
    是任何 \(n-r\) 个数.
    我们考虑%
    由 \(r\)~个 \(r\)~元 \({\leq} 1\)~次方程作成的方程组
    \begin{align*}
        \sum_{\ell = 1}^{r}
        {[A]_{p,j_\ell} y_\ell}
        = [B]_{p,1}
        - \sum_{r < \ell \leq n}
        {[A]_{p,j_\ell} f_{j_\ell}}.
    \end{align*}
    利用阵等式, 我们可写
    \begin{align*}
        A_r
        \begin{bmatrix}
            y_1    \\
            y_2    \\
            \vdots \\
            y_r    \\
        \end{bmatrix}
        =
        \begin{bmatrix}
            z_{i_1} \\
            z_{i_2} \\
            \vdots  \\
            z_{i_r} \\
        \end{bmatrix},
    \end{align*}
    其中,
    \begin{align*}
        z_p = [B]_{p,1}
        - \sum_{r < \ell \leq n}
        {[A]_{p,j_\ell} f_{j_\ell}}.
    \end{align*}
    因为 \(\det {(A_r)} \neq 0\),
    故, 由 Cramer 公式,
    存在 \(r\) 个数
    \(f_{j_1}\), \(f_{j_2}\), \(\dots\), \(f_{j_r}\),
    使
    \begin{align*}
        \sum_{\ell = 1}^{r}
        {[A]_{p,j_\ell} f_{j_\ell}}
            = [B]_{p,1}
        - \sum_{r < \ell \leq n}
        {[A]_{p,j_\ell} f_{j_\ell}}.
    \end{align*}
    从而
    \begin{align*}
        \sum_{\ell = 1}^{n}
        {[A]_{p,j_\ell} f_{j_\ell}}
            = [B]_{p,1}.
    \end{align*}

    设 \(n\) 个数
    \(c_1\), \(c_2\), \(\dots\), \(c_n\)
    适合
    \begin{align*}
        [A]_{i_1,1} c_1 + [A]_{i_1,2} c_2 + \dots + [A]_{i_1,n} c_n & = [B]_{i_1,1}, \\
        [A]_{i_2,1} c_1 + [A]_{i_2,2} c_2 + \dots + [A]_{i_2,n} c_n & = [B]_{i_2,1}, \\
                                                                    & \dots,         \\
        [A]_{i_r,1} c_1 + [A]_{i_r,2} c_2 + \dots + [A]_{i_r,n} c_n & = [B]_{i_r,1}.
    \end{align*}
    作 \(n \times 1\)~阵 \(C\),
    其中, \([C]_{i,1} = c_i\).
    我们证明, \(AC = B\).

    若 \(i\) 等于某 \(i_s\), 则
    \begin{align*}
        [AC]_{i,1}
        = {} &
        \sum_{\ell = 1}^{n}
        {[A]_{i,\ell} [C]_{\ell,1}}
        \\
        = {} &
        \sum_{\ell = 1}^{n}
        {[A]_{i,\ell} c_\ell}
        \\
        = {} & [B]_{i,1}.
    \end{align*}
    若 \(i\) 不等于任何 \(i_s\), 则
    \begin{align*}
        [AC]_{i,1}
        = {} &
        \sum_{\ell = 1}^{n}
        {[A]_{i,\ell} [C]_{\ell,1}}
        \\
        = {} &
        \sum_{\ell = 1}^{n}
        {[G]_{i,\ell} [C]_{\ell,1}}
        \\
        = {} &
        \sum_{\ell = 1}^{n}
        {
        \left( \sum_{s = 1}^{r}
        {
            [G]_{i_s,\ell} d_{i,s}
        } \right) [C]_{\ell,1}
        }
        \\
        = {} &
        \sum_{\ell = 1}^{n}
        {
        \sum_{s = 1}^{r}
        {
        [G]_{i_s,\ell} d_{i,s} [C]_{\ell,1}
        }
        }
        \\
        = {} &
        \sum_{s = 1}^{r}
        {
        \sum_{\ell = 1}^{n}
        {
        [G]_{i_s,\ell} d_{i,s} [C]_{\ell,1}
        }
        }
        \\
        = {} &
        \sum_{s = 1}^{r}
        {
        \sum_{\ell = 1}^{n}
        {
        [A]_{i_s,\ell} d_{i,s} c_\ell
        }
        }
        \\
        = {} &
        \sum_{s = 1}^{r}
        {
        \sum_{\ell = 1}^{n}
        {
        [A]_{i_s,\ell} c_\ell d_{i,s}
        }
        }
        \\
        = {} &
        \sum_{s = 1}^{r}
        {
        \left(\sum_{\ell = 1}^{n}
        {
            [A]_{i_s,\ell} c_\ell
        }\right)
        d_{i,s}
        }
        \\
        = {} &
        \sum_{s = 1}^{r}
        {[B]_{i_s,1} d_{i,s}}
        \\
        = {} &
        \sum_{s = 1}^{r}
        {[G]_{i_s,n+1} d_{i,s}}
        \\
        = {} &
        [G]_{i,n+1}
        \\
        = {} &
        [B]_{i,1}.
        \qedhere
    \end{align*}
\end{proof}

结合前几节的讨论, 我们可以得到判断
\(AX = B\) 是否有解,
与有解时它的解是否唯一的方法
(定性的理论):

% restatable
\begin{restatable}{theorem}%
    {TheoremQualitativeTheoryOfLinearSystemOfEquations}
    设 \(A\) 是 \(m \times n\)~阵.
    设 \(B\) 是 \(m \times 1\)~阵.
    作一个 \(m \times (n+1)\)~阵 \(G\),
    其中,
    \begin{align*}
        [G]_{i,j}
        = \begin{cases}
              [A]_{i,j}, & j \leq n;  \\
              [B]_{i,1}, & j = n + 1.
          \end{cases}
    \end{align*}
    (通俗地, 在 \(A\)~的最后一列的右侧加入一列 \(B\),
    得到尺寸较大的阵 \(G\).)
    设
    \(A\) 有一个行列式非零的 \(r\)~级子阵,
    但没有行列式非零的 \(r+1\)~级子阵.

    (1)
    若 \(G\)~也没有行列式非零的 \(r+1\)~级子阵,
    则 \(AX = B\) 有解.
    进一步地, 若 \(r < n\),
    则 \(AX = B\) 的解不唯一;
    若 \(r = n\),
    则 \(AX = B\) 的解唯一.

    (2)
    若 \(G\)~有一个行列式非零的 \(r+1\)~级子阵,
    则 \(AX = B\) 无解.
\end{restatable}

\KunAsteriskoEnEnhavtabelo
\section{\texorpdfstring{由 \(m\)~个 \(n\)~元
      \({\leq} 1\)~次方程作成的方程组 (4)}%
  {由 m 个 n 元 ≤1 次方程作成的方程组 (4)}}
\SenAsteriskoEnEnhavtabelo

\maldevigalegajxo

前面, 我们得到了线性方程组的解的定性的理论:

\TheoremQualitativeTheoryOfLinearSystemOfEquations*

本节, 我们作定量的讨论:
当 \(AX = B\) 有解时,
我们试作出公式, 以表示它的解.

设 \(A\) 是一个 \(m \times n\)~阵.
设 \(B\) 是一个 \(m \times 1\)~阵.
设 \(AX = B\) 有解.

若 \(A = 0\),
则因 \(AX = B\) 有解,
必有 \(B = 0\).
此时, 显然, 每个 \(n \times 1\)~阵都是解.
下设 \(A \neq 0\).
设 \(A\) 有一个行列式非零的 \(r\)~级子阵
\begin{align*}
    T = A\binom{i_1,\dots,i_r}{j_1,\dots,j_r}
\end{align*}
(其中,
\(1 \leq i_1 < \dots < i_r \leq m\),
且 \(1 \leq j_1 < \dots < j_r \leq n\)),
但没有行列式非零的 \(r+1\)~级子阵.
为方便, 我们作一个 \(r \times n\)~阵
\begin{align*}
    U = A\binom{i_1,\dots,i_r}{1,\dots,n},
\end{align*}
与一个 \(r \times 1\)~阵
\begin{align*}
    V = B\binom{i_1,\dots,i_r}{1}
    % =
    % \begin{bmatrix}
    %     [B]_{i_1,1} \\
    %     [B]_{i_2,1} \\
    %     \vdots      \\
    %     [B]_{i_r,1} \\
    % \end{bmatrix}
    .
\end{align*}
注意, \(U\) 有一个行列式非零的 \(r\)~级子阵 \(T\).

我们回想上节的讨论.
为解方程组 \(AX = B\), 即
\begin{equation*}
    \left\{
    \begin{aligned}
        [A]_{1,1} x_1 + [A]_{1,2} x_2 + \dots + [A]_{1,n} x_n & = [B]_{1,1}, \\
        [A]_{2,1} x_1 + [A]_{2,2} x_2 + \dots + [A]_{2,n} x_n & = [B]_{2,1}, \\
                                                              & \dots,       \\
        [A]_{m,1} x_1 + [A]_{m,2} x_2 + \dots + [A]_{m,n} x_n & = [B]_{m,1},
    \end{aligned}
    \right.
\end{equation*}
我们考虑了由这 \(m\)~个方程的%
第~\(i_1\), \(i_2\), \(\dots\), \(i_r\)~个方程%
作成的方程组
\begin{equation*}
    \left\{
    \begin{aligned}
        [A]_{i_1,1} x_1 + [A]_{i_1,2} x_2 + \dots + [A]_{i_1,n} x_n & = [B]_{i_1,1}, \\
        [A]_{i_2,1} x_1 + [A]_{i_2,2} x_2 + \dots + [A]_{i_2,n} x_n & = [B]_{i_2,1}, \\
                                                                    & \dots,         \\
        [A]_{i_r,1} x_1 + [A]_{i_r,2} x_2 + \dots + [A]_{i_r,n} x_n & = [B]_{i_r,1},
    \end{aligned}
    \right.
\end{equation*}
即 \(UX = V\).
当时, 我们已证明,
\(UX = V\) 的解都是 \(AX = B\) 的解.
反过来, \(AX = B\) 的解显然都是 \(UX = V\) 的解,
因为后者的方程全部都是来自前者的.
这么看来, \(AX = B\) 跟 \(UX = V\) 有相同的解.
所以, 研究 \(AX = B\) 的解的公式,
相当于研究 \(UX = V\) 的解的公式.

若 \(r = n\),
则我们直接用 Cramer 公式,
即可写出 \(UX = V\) 的唯一的解
(注意, 此时 \(U\) 是一个 \(r \times r\)~阵)
\begin{align*}
    X = (\det {(U)})^{-1} \operatorname{adj} {(U)}\,V.
\end{align*}
我们也可较直接地表示此解.
设 \(U\{k,V\}\) 是以 \(r \times 1\)~阵 \(V\)
代 \(U\) 的列~\(k\) 后得到的阵.
则
\begin{align*}
    x_k = \frac{\det {(U\{k,V\})}}{\det {(U)}}.
\end{align*}

下设 \(r < n\).
我们试写出方程组 \(UX = V\) 的所有的解.

\begin{theorem}
    设 \(V\) 是一个 \(r \times 1\)~阵.
    设 \(U\) 是一个 \(r \times n\)~阵,
    其中, \(r < n\),
    且 \(U\) 有一个行列式非零的 \(r\)~级子阵
    \begin{align*}
        T = U\binom{1,\dots,r}{j_1,\dots,j_r},
    \end{align*}
    其中, \(1 \leq j_1 < \dots < j_r \leq n\).
    从 \(1\), \(2\), \(\dots\), \(n\)
    去除 \(j_1\), \(j_2\), \(\dots\), \(j_r\)
    后, 还剩 \(n - r\)~个数.
    我们从小到大地叫这 \(n - r\)~个数为
    \(j_{r+1}\), \(\dots\), \(j_n\).
    再设 \(U\)~的%
    列~\(1\), \(2\), \(\dots\), \(n\)
    分别是 \(u_1\), \(u_2\), \(\dots\), \(u_n\).

    (1)
    设 \(c_{j_{r+1}}\), \(\dots\), \(c_{j_n}\)
    是 \(n-r\)~个常数.
    作 \(n \times 1\)~阵 \(C\), 其中,
    \begin{align*}
        [C]_{j_k,1}
        = \begin{dcases}
              \frac{\det {(T\{k,V\})}}{\det {(T)}}
              + \sum_{r < \ell \leq n}
              {c_{j_\ell}
              \frac{-\det {(T\{k,u_{j_\ell}\})}}{\det {(T)}}},
               & k \leq r; \\
              c_{j_k},
               & k > r.
          \end{dcases}
    \end{align*}
    其中, \(T\{k,Y\}\) 是以 \(r \times 1\)~阵 \(Y\)
    代 \(T\) 的列~\(k\) 后得到的阵.
    则 \(UC = V\).

    (2)
    若 \(n \times 1\)~阵 \(D\) 适合 \(UD = V\),
    则存在 \(n-r\)~个数
    \(c_{j_{r+1}}\), \(\dots\), \(c_{j_n}\),
    使
    \begin{align*}
        [D]_{j_k,1}
        = \begin{dcases}
              \frac{\det {(T\{k,V\})}}{\det {(T)}}
              + \sum_{r < \ell \leq n}
              {c_{j_\ell}
              \frac{-\det {(T\{k,u_{j_\ell}\})}}{\det {(T)}}},
               & k \leq r; \\
              c_{j_k},
               & k > r.
          \end{dcases}
    \end{align*}
    或者, \(UX = V\) 的每个解%
    都可被 (1) 中的公式表示.
\end{theorem}

\begin{proof}
    (1)
    我们可改写方程组 \(UX = V\) 为
    \begin{align*}
        \sum_{\ell = 1}^{r}
        {[U]_{p,j_\ell} x_{j_\ell}}
            = [V]_{p,1}
        - \sum_{r < \ell \leq n}
        {[U]_{p,j_\ell} x_{j_\ell}},
    \end{align*}
    其中, \(p = 1\), \(2\), \(\dots\), \(r\),
    下同.
    考虑%
    由 \(r\)~个 \(r\)~元 \({\leq} 1\)~次方程作成的方程组
    \begin{align*}
        \sum_{\ell = 1}^{r}
        {[U]_{p,j_\ell} y_\ell}
        = [V]_{p,1}
        - \sum_{r < \ell \leq n}
        {[U]_{p,j_\ell} c_{j_\ell}}.
    \end{align*}
    利用阵等式, 我们可写
    \begin{align*}
        T
        \begin{bmatrix}
            y_1    \\
            y_2    \\
            \vdots \\
            y_r    \\
        \end{bmatrix}
        =
        V -
        \sum_{r < \ell \leq n}
        {c_{j_\ell} u_{j_\ell}}.
    \end{align*}
    因为 \(\det {(T)} \neq 0\),
    故, 由 Cramer 公式,
    存在 \(r\) 个数
    \(c_{j_1}\), \(c_{j_2}\), \(\dots\), \(c_{j_r}\),
    使
    \begin{align*}
        \sum_{\ell = 1}^{r}
        {[U]_{p,j_\ell} c_{j_\ell}}
            = [V]_{p,1}
        - \sum_{r < \ell \leq n}
        {[U]_{p,j_\ell} c_{j_\ell}},
    \end{align*}
    其中,
    \begin{align*}
        c_{j_k}
        = {} &
        (\det {(T)})^{-1}
        \det {
            \left(
            T\left\{
            k,
            V -
            \sum_{r < \ell \leq n}
            {c_{j_\ell} u_{j_\ell}}
            \right\}
            \right)
        }
        \\
        = {} &
        (\det {(T)})^{-1}
        \left(
        \det {(T\{k,V\})}
        -
        \sum_{r < \ell \leq n}
        {c_{j_\ell}
        \det {(T\{k,u_{j_\ell}\})}
        }
        \right)
        \\
        = {} &
        \frac{\det {(T\{k,V\})}}{\det {(T)}}
        + \sum_{r < \ell \leq n}
        {c_{j_\ell}
        \frac{-\det {(T\{k,u_{j_\ell}\})}}{\det {(T)}}}.
    \end{align*}
    从而
    \begin{align*}
        \sum_{\ell = 1}^{n}
        {[U]_{p,j_\ell} c_{j_\ell}}
            = [V]_{p,1}.
    \end{align*}
    注意, \([C]_{j_k,1} = c_{j_k}\).
    所以, \(UC = V\).

    (2)
    设 \(D\) 适合 \(UD = V\).
    则
    \begin{align*}
        \sum_{\ell = 1}^{n}
        {[U]_{p,j_\ell} [D]_{j_\ell,1}}
            = [V]_{p,1}.
    \end{align*}
    从而
    \begin{align*}
        \sum_{\ell = 1}^{r}
        {[U]_{p,j_\ell} [D]_{j_\ell,1}}
            = [V]_{p,1}
        - \sum_{r < \ell \leq n}
        {[U]_{p,j_\ell} [D]_{j_\ell,1}}.
    \end{align*}

    取 \(c_{j_\ell} = [D]_{j_\ell,1}\),
    \(\ell > r\).
    考虑方程组
    \begin{align*}
        \sum_{\ell = 1}^{r}
        {[U]_{p,j_\ell} y_\ell}
        = [V]_{p,1}
        - \sum_{r < \ell \leq n}
        {[U]_{p,j_\ell} c_{j_\ell}}.
    \end{align*}
    由 (1), 我们知道
    \begin{align*}
        y_k =
        \frac{\det {(T\{k,V\})}}{\det {(T)}}
        + \sum_{r < \ell \leq n}
        {c_{j_\ell}
        \frac{-\det {(T\{k,u_{j_\ell}\})}}{\det {(T)}}}
    \end{align*}
    (其中, \(k = 1\), \(2\), \(\dots\), \(r\), 下同)
    是一个解;
    另一方面,
    \(y_k = [D]_{j_k,1}\) 也是一个解.
    因为 \(\det {(T)} \neq 0\),
    故, 由 Cramer 公式,
    这二个解应是相同的,
    即
    \begin{equation*}
        [D]_{j_k,1} =
        \frac{\det {(T\{k,V\})}}{\det {(T)}}
        + \sum_{r < \ell \leq n}
        {c_{j_\ell}
        \frac{-\det {(T\{k,u_{j_\ell}\})}}{\det {(T)}}}.
        \qedhere
    \end{equation*}
\end{proof}

注意, 当 \(r = n\) 时,
不可能有数 \(\ell\) 适合 \(r < \ell \leq n\).
既然没有数被加, 那么, 形如
\begin{align*}
    \sum_{r < \ell \leq n} {f(\ell)}
\end{align*}
的式是 \(0\).
用这个约定,
我们可统一地写%
在 \(r = n\) 与 \(r < n\) 这二个情形下%
解的公式.

\begin{theorem}
    设 \(V\) 是一个 \(r \times 1\)~阵.
    设 \(U\) 是一个 \(r \times n\)~阵,
    其中, \(r \leq n\),
    且 \(U\) 有一个行列式非零的 \(r\)~级子阵
    \begin{align*}
        T = U\binom{1,\dots,r}{j_1,\dots,j_r},
    \end{align*}
    其中, \(1 \leq j_1 < \dots < j_r \leq n\).
    从 \(1\), \(2\), \(\dots\), \(n\)
    去除 \(j_1\), \(j_2\), \(\dots\), \(j_r\)
    后, 还剩 \(n - r\)~个数.
    我们从小到大地叫这 \(n - r\)~个数为
    \(j_{r+1}\), \(\dots\), \(j_n\).
    再设 \(U\)~的%
    列~\(1\), \(2\), \(\dots\), \(n\)
    分别是 \(u_1\), \(u_2\), \(\dots\), \(u_n\).
    那么, \(UX = V\) 的解恰为
    \begin{align*}
        x_{j_k}
        = \begin{dcases}
              \frac{\det {(T\{k,V\})}}{\det {(T)}}
              + \sum_{r < \ell \leq n}
              {c_{j_\ell}
              \frac{-\det {(T\{k,u_{j_\ell}\})}}{\det {(T)}}},
               & k \leq r; \\
              c_{j_k},
               & k > r,
          \end{dcases}
    \end{align*}
    其中, \(c_{j_{r+1}}\), \(\dots\), \(c_{j_n}\)
    是任何 \(n-r\) 个数.
\end{theorem}

最后, 我们总结这几节的关于线性方程组的主要结果.

\begin{theorem}
    设 \(A\) 是 \(m \times n\)~阵.
    设 \(B\) 是 \(m \times 1\)~阵.
    作一个 \(m \times (n+1)\)~阵 \(G\),
    其中,
    \begin{align*}
        [G]_{i,j}
        = \begin{cases}
              [A]_{i,j}, & j \leq n;  \\
              [B]_{i,1}, & j = n + 1.
          \end{cases}
    \end{align*}
    (通俗地, 在 \(A\)~的最后一列的右侧加入一列 \(B\),
    得到尺寸较大的阵 \(G\).)
    设 \(A\) 有一个行列式非零的 \(r\)~级子阵,
    但 \(A\) 没有行列式非零的 \(r+1\)~级子阵;
    设 \(G\) 有一个行列式非零的 \(s\)~级子阵,
    但 \(G\) 没有行列式非零的 \(s+1\)~级子阵.
    设 \(X\) 是未知的 \(n \times 1\)~阵.

    (1)
    因为 \(A\) 是 \(G\) 的子阵,
    故 \(A\) 的子阵也是 \(G\) 的子阵.
    则 \(s = r\) 或 \(s > r\).
    % 因为 \(G\) 恰比 \(A\) 多一列,
    % 故 \(s \leq r + 1\)
    % (若 \(G\) 有一个行列式非零的 \(r + 2\)~级子阵,
    % 则 \(A\) 应有一个行列式非零的 \(r + 1\)~级子阵;
    % 这是矛盾).

    (2)
    若 \(s > r\), 则 \(AX = B\) 无解;
    也就是, 若 \(AX = B\) 有解, 必 \(s = r\).

    (3)
    若 \(s = r\), 则 \(AX = B\) 有解;
    也就是, 若 \(AX = B\) 无解, 必 \(s > r\).

    (4)
    设 \(s = r\).
    则 \(AX = B\) 有解.
    若 \(r = n\), 则 \(AX = B\) 的解唯一;
    若 \(r < n\), 则 \(AX = B\) 的解不唯一.

    (5)
    设 \(A \neq 0\).
    设 \(A\) 有一个行列式非零的 \(r\)~级子阵
    \begin{align*}
        T = A\binom{i_1,\dots,i_r}{j_1,\dots,j_r}
    \end{align*}
    (其中, \(1 \leq i_1 < \dots < i_r \leq m\),
    且 \(1 \leq j_1 < \dots < j_r \leq n\)),
    但 \(G\) 没有行列式非零的 \(r+1\)~级子阵.
    则 \(s = r\).
    则 \(AX = B\) 有解.

    记 \(r \times n\)~阵
    \begin{align*}
        U = A\binom{i_1,\dots,i_r}{1,\dots,n}.
    \end{align*}
    于是, 可视 \(T\) 为 \(U\) 的一个%
    行列式非零的 \(r\)~级子阵.
    再记 \(r \times 1\)~阵
    \begin{align*}
        V = B\binom{i_1,\dots,i_r}{1}
        % =
        % \begin{bmatrix}
        %     [B]_{i_1,1} \\
        %     [B]_{i_2,1} \\
        %     \vdots      \\
        %     [B]_{i_r,1} \\
        % \end{bmatrix}
        .
    \end{align*}
    则 \(UX = V\) 有解,
    \(AX = B\) 的解是 \(UX = V\) 的解,
    且 \(UX = V\) 的解是 \(AX = B\) 的解.

    从 \(1\), \(2\), \(\dots\), \(n\)
    去除 \(j_1\), \(j_2\), \(\dots\), \(j_r\)
    后, 还剩 \(n - r\)~个数.
    我们从小到大地叫这 \(n - r\)~个数为
    \(j_{r+1}\), \(\dots\), \(j_n\).
    再设 \(U\)~的%
    列~\(1\), \(2\), \(\dots\), \(n\)
    分别是 \(u_1\), \(u_2\), \(\dots\), \(u_n\).
    那么, \(UX = V\) 的解 (即 \(AX = B\) 的解) 恰为
    \begin{align*}
        x_{j_k}
        = \begin{dcases}
              \frac{\det {(T\{k,V\})}}{\det {(T)}}
              + \sum_{r < \ell \leq n}
              {c_{j_\ell}
              \frac{-\det {(T\{k,u_{j_\ell}\})}}{\det {(T)}}},
               & k \leq r; \\
              c_{j_k},
               & k > r,
          \end{dcases}
    \end{align*}
    其中,
    \(c_{j_{r+1}}\), \(\dots\), \(c_{j_n}\)
    是任何 \(n-r\) 个数,
    \(T\{k,Y\}\) 是以 \(r \times 1\)~阵 \(Y\)
    代 \(T\) 的列~\(k\) 后得到的阵,
    且 \(x_i = [X]_{i,1}\),
    \(i = 1\), \(2\), \(\dots\), \(n\).
\end{theorem}

\KunAsteriskoEnEnhavtabelo
\section{\texorpdfstring{由 \(m\)~个 \(n\)~元
      \({\leq} 1\)~次方程作成的方程组 (5)}%
  {由 m 个 n 元 ≤1 次方程作成的方程组 (5)}}
\SenAsteriskoEnEnhavtabelo

\maldevigalegajxo

线性方程组的故事要结束了.
我们已解决线性方程组的三个重要的问题:

(1)
一个线性方程组何时有解?

(2)
一个线性方程组有解时, 它的解是否唯一?

(3)
一个线性方程组有解时, 我们如何找到它的全部的解?

我们用行列式回答了这三个问题,
作了线性方程组的一个理论.
所以, 理论地, 给了一个线性方程组,
我们可以计算一些行列式以确定它是否有解,
且有解时它的解为何.

不过,
计算一个方阵的行列式并不是什么简单的事情:
\(3\)~级阵的行列式的具体的公式含 \(6\)~项,
且 \(4\)~级阵的行列式的具体的公式含 \(24\)~项.

不过, 此理论仍然是重要的.
或许, 您还记得,
任给一个阵~\(A\),
必存在唯一的非负整数 \(r\),
使 \(A\)~有一个行列式非零的 \(r\)~级子阵,
但 \(A\)~没有行列式非零的 \(r+1\)~级子阵.
聪明的数学人意识到,
此 \(r\) 是重要的,
故特别地为它起了一个名字, ``秩''.
然后, 他们研究了秩,
发现:
(a)
可以施适当的变换于阵~\(A\),
得到一个新阵~\(K\),
且 \(A\) 与 \(K\) 有相同的秩;
(b)
``适当的变换'' 跟计算方阵的行列式比,
有较少的计算量
(具体地, 加、减、乘、除的次数);
(c)
\(K\)~的秩不难计算,
且可被直接看出.
于是, 联合这个发现与我们的理论,
就是一个更好的线性方程组的理论.
(``适当的变换'' 也是解线性方程组的好的方法.)
若您想知道更多,
您可以见线性代数教材.

我就说这么多.

\end{document}

This material discusses
systems of m linear equations with n unknowns.
As is indicated in the notice of dependence,
one can directly learn this material
after knowing what a system of linear equations is
and Cramer's formula.
Discussing
systems of m linear equations with n unknowns
mainly with determinants
is of theoretical interest.
I do not explicitly define what the rank of a matrix is
in this material,
but readers who are acquainted
with the basics of linear algebra
should be able to realise
what role rank plays in the discussion.
